\documentclass[11pt]{article}
\usepackage[T1]{fontenc}
\usepackage[utf8]{inputenc}
\usepackage{lmodern}
\usepackage[a4paper,margin=1in]{geometry}
\usepackage{microtype}
\usepackage{amsmath,amssymb,amsthm,mathtools}
\usepackage{booktabs,array}
\usepackage[dvipsnames]{xcolor}
\usepackage[colorlinks=true,linkcolor=MidnightBlue,citecolor=MidnightBlue,urlcolor=MidnightBlue]{hyperref}
\usepackage{enumitem}
\usepackage{setspace}
\usepackage{fancyhdr}
\usepackage{titlesec}
\setlength{\headheight}{14pt}
\pagestyle{fancy}
\fancyhf{}
\fancyhead[L]{Repaired split-forest proof: items 4--8 of round 5}
\fancyhead[R]{May 2026}
\fancyfoot[C]{\thepage}
\titleformat{\section}{\large\bfseries\color{MidnightBlue}}{\thesection}{0.6em}{}
\titleformat{\subsection}{\normalsize\bfseries\color{MidnightBlue}}{\thesubsection}{0.5em}{}
\setlist[itemize]{topsep=4pt,itemsep=2pt,leftmargin=1.5em}
\setlist[enumerate]{topsep=4pt,itemsep=2pt,leftmargin=1.8em}
\onehalfspacing

\newtheorem{theorem}{Theorem}[section]
\newtheorem{proposition}[theorem]{Proposition}
\newtheorem{lemma}[theorem]{Lemma}
\newtheorem{corollary}[theorem]{Corollary}
\theoremstyle{definition}
\newtheorem{definition}[theorem]{Definition}
\newtheorem{construction}[theorem]{Construction}
\theoremstyle{remark}
\newtheorem{remark}[theorem]{Remark}
\newtheorem{example}[theorem]{Example}

\newcommand{\DR}{\mathsf{DR}}
\newcommand{\PO}{\mathsf{PO}}
\newcommand{\PP}{\mathsf{PP}}
\newcommand{\PPI}{\mathsf{PPI}}
\newcommand{\EQ}{\mathsf{EQ}}
\newcommand{\Base}{\mathsf{B}}
\newcommand{\Inv}{\operatorname{Inv}}
\newcommand{\comp}{\operatorname{comp}}
\newcommand{\ALCIRCC}[1]{\mathcal{ALCI}_{\mathsf{RCC#1}}}
\newcommand{\Vrt}{\mathsf{Vrt}}
\newcommand{\Sib}{\mathsf{Sib}}
\newcommand{\Front}{\mathsf{Front}}
\newcommand{\Mos}{\mathsf{Mos}}
\newcommand{\Ctx}{\mathsf{Ctx}}
\newcommand{\Pos}{\mathsf{Pos}}
\newcommand{\Trip}{\mathsf{Trip}}
\newcommand{\Pair}{\mathsf{Pair}}
\newcommand{\Port}{\mathsf{Port}}
\newcommand{\Mate}{\mathsf{Mate}}
\newcommand{\Reach}{\mathsf{Reach}}
\newcommand{\Side}{\mathsf{Side}}
\newcommand{\Q}{\mathcal{Q}}
\newcommand{\U}{\mathcal{U}}
\newcommand{\T}{\mathcal{T}}
\newcommand{\md}{\operatorname{md}}
\newcommand{\Cl}{\operatorname{cl}}
\newcommand{\Typ}{\operatorname{Typ}}
\newcommand{\Pend}{\mathsf{Pend}}
\newcommand{\chk}{\mathsf{chk}}
\newcommand{\eqc}{\equiv_{\EQ}}
\newcommand{\eqQEQ}{\equiv^{\Q}_{\EQ}}
\newcommand{\extProf}{\widehat\Pi}
\newcommand{\extSigma}{\widehat\Sigma}
\newcommand{\lab}{\operatorname{lab}}
\newcommand{\Stratum}{\mathsf{Strat}}

\begin{document}

\begin{center}
{\LARGE\bfseries Repaired Split-Forest Decidability Proof for $\ALCIRCC{5}$:\\[2pt]
Items 4--8 of the Round-5 Review\par}
\vspace{0.5em}
{\large A delta to the 31-page round-4 proof~\cite{GPT55Repaired},\\
companion to the (M6$'$) v2 delta~\cite{V2Delta}\par}
\vspace{0.75em}
{\normalsize Claude (Anthropic, Opus~4.7), prompted by Michael Wessel}\par
\vspace{0.45em}
{\small May~2026}
\end{center}

\begin{center}
\fbox{\parbox{0.92\textwidth}{%
\textbf{Disclaimer.}\;
This document was drafted by Claude (Anthropic), prompted by Michael
Wessel.  The results have not been independently verified by human
domain experts.  We invite scrutiny and corrections.}}
\end{center}

\begin{abstract}
The round-5 review of the 31-page repaired split-forest decidability
proof~\cite{GPT55Repaired} for $\ALCIRCC{5}$~\cite{GPT55Round5Review}
lists eight recommended repairs.  Item~1 (TeX/PDF mismatch) is
editorial and already self-consistent in the repository; items~2, 3
and~9 are addressed by the (M6$'$) recursive-verticalization v2
delta~\cite{V2Delta}.  This delta paper closes the remaining
mathematical items 4--8:
\begin{enumerate}[label=(\arabic*),leftmargin=2em]
\item[\textbf{4.}] \emph{Non-circular mosaic patchwork.}  We replace
the implicit ``support against $\T_\Q$'' clause~(SC4) of round 4 by
an explicit two-step construction: first, a finite abstract
pair-label function $L_\Q:\Pos_\Q^2\to\Base$ is declared as
certificate data; second, an explicit finite validity clause~(V11)
requires every abstract triple in $\Trip_\Q$ to satisfy the RCC5
composition table.  The mosaic patchwork theorem is then a
non-circular consequence.
\item[\textbf{5.}] \emph{Strengthened overlap amalgamation.}  We
strengthen hypothesis~(O3) of round 4 to require a single
cross-label assignment $L_{12}:(P_1\setminus P_0)\times(P_2\setminus
P_0)\to\Base$ compatible with all overlap nodes and all mixed
triples simultaneously.  This becomes an additional explicit
validity clause~(V12) on certificates.
\item[\textbf{6.}] \emph{Joint equality validity clause.}  We
promote round-4 assumption~(A4) of the typed-$\EQ$ quotient theorem
to an explicit finite validity clause~(V13): for every equality
class $E$ of ports and every $\exists R.D$ common to the type of
$E$'s members, there is a declared mate/target pair that works
uniformly for all ports in $E$.
\item[\textbf{7.}] \emph{Half-open cycle convention.}  We rewrite
the request-fulfilled repetition lemma with the half-open
$[i,j)$ convention.  The pumped path is $u_0,\ldots,u_{i-1},
(u_i,\ldots,u_{j-1})^\omega$, and $u_j$ is identified with the next
lap's copy of $u_i$ at the profile level only.  The endpoint case
is handled explicitly.
\item[\textbf{8.}] \emph{Global finite regularization.}  We add a
global representative-selection lemma: for the finitely many
$(\text{profile},\text{request-status})$ classes
$\extProf\in\extSigma$, one representative finite context and one
representative request-fulfilled cycle suffices, simultaneously,
for every infinite vertical ray of the split forest.  The proof is
a direct pigeonhole over $\extSigma$, not via $\omega$-automata.
\end{enumerate}
None of items 4--8 affects~(M6$'$); the architecture of the v2
delta paper carries over, and the validity-clause numbering is
extended from (V1)--(V10) to (V1)--(V13).  The decision procedure
remains a finite enumeration of finite-graph-checkable certificates;
$B(C_0)\leq 2^{2^{p(n)}}$ for a computable polynomial~$p$ is
preserved.
\end{abstract}

\tableofcontents
\bigskip

%% ==================================================================
\section{Scope and notation}
\label{sec:scope}
%% ==================================================================

This document is a delta to the 31-page round-4 paper
\cite{GPT55Repaired} and the (M6$'$) v2 delta paper \cite{V2Delta}.
We use the round-4 notation throughout: $\Q$ is a finite generated
split-forest certificate; $\Mos_\Q$ is its support-closed mosaic
family; $\T_\Q=(\Pos_\Q,\Pair_\Q,\Trip_\Q)$ is the abstract triple
graph; $\U(\Q)$ is the unfolding; $\eqc^\Q$ is the explicit equality
congruence on ports of $\Q$.  All other notation is as in
\cite{GPT55Repaired}.

Where the v2 delta paper restated a definition or a clause of
\cite{GPT55Repaired}, we work with the v2 restated form
(e.g., recursive verticalization (M6$'$) replaces (M6); validity
clauses (V3$'$), (V4$'$), (V5$'$), (V9.Eb$'$) replace (V3), (V4),
(V5), (V9.Eb)).

%% ==================================================================
\section{Item 4: non-circular mosaic patchwork}
\label{sec:item4}
%% ==================================================================

\subsection{The circularity diagnosed by round 5}

Round-4 \cite{GPT55Repaired} proves the support-closed mosaic
patchwork theorem (Theorem~5.8) in two steps:
\begin{enumerate}
\item the mosaic relational layer satisfies the Renz--Nebel
hypotheses (RN1)--(RN3) on every mosaic (Lemma~5.5);
\item the overlap-amalgamation lemma (Lemma~5.6) combines two
mosaics over a shared boundary;
\item the support-closure clause (SC4) requires that every abstract
triple in $\Trip_\Q$ is represented by some mosaic $M_T\in\Mos$ with
matching pair labels (Definition~5.4).
\end{enumerate}

The round-5 review~\cite{GPT55Round5Review} objects that the
patchwork theorem uses (SC4) to argue that concrete unfolding
triples are RCC5-composition-consistent, while (SC4) itself
\emph{presupposes} that the mosaic $M_T$ assigning labels to an
abstract triple realises composition-consistent labels.  The
labelling on $\T_\Q$ is implicit, contained in the mosaic data
itself rather than declared as an independent function.  The risk
is that the patchwork claim is established by appeal to mosaic
composition consistency, which is in turn what the patchwork claim
asserts at the abstract level.

\subsection{The repair: explicit $L_\Q$ and a composition-table validity clause}

We declare an explicit \emph{abstract pair-label function} on the
abstract triple graph, separately from any mosaic, and require it
to satisfy the composition table.  The certificate's pair labels
become primary data; mosaics become local realisations of those
labels.

\begin{definition}[Abstract pair-label function]\label{def:Lq}
The certificate $\Q$ declares an explicit function
\[
  L_\Q:\Pos_\Q^2\;\longrightarrow\;\Base
\]
on pairs of abstract position symbols, satisfying:
\begin{enumerate}[label=(L\arabic*)]
\item $L_\Q(\pi,\pi)=\EQ$;
\item $L_\Q(\pi',\pi)=\Inv(L_\Q(\pi,\pi'))$;
\item if $\pi,\pi'$ are linked by an explicit equality-port chain in
$E_\Q$, then $L_\Q(\pi,\pi')=\EQ$.
\end{enumerate}
\end{definition}

\begin{definition}[Compatibility with mosaics]\label{def:LM-Lq}
For every mosaic $M=(P,\tau,L,\mathcal V)\in\Mos_\Q$ and every pair
$p,q\in P$ realising abstract position symbols
$\pi_p,\pi_q\in\Pos_\Q$:
\[
  L_M(p,q)\;=\;L_\Q(\pi_p,\pi_q).
\]
That is, mosaic labels are restrictions of the global $L_\Q$ to the
mosaic's position set.
\end{definition}

The labels in $L_\Q$ are part of the certificate $\Q$ and are part
of the finite data the validity check enumerates.  We now require
explicit composition-table checking on $\T_\Q$ as a finite
validity clause:

\paragraph{(V11) Composition table on $\T_\Q$.}
For every ordered triple $(\pi_1,\pi_2,\pi_3)\in\Pos_\Q^3$,
\[
  L_\Q(\pi_1,\pi_3)\;\in\;\comp\bigl(L_\Q(\pi_1,\pi_2),
                                       L_\Q(\pi_2,\pi_3)\bigr).
\]
This is a finite quantifier over $|\Pos_\Q|^3\leq 2|S|\cdot(n+1)
\cdot|\Ctx_\Q|\cdot k_{\rm strat}\cdot |\Pos_\Q|^2$
abstract triples, each requiring a constant-time table lookup.

\begin{lemma}[(V11) implies (M3) on every mosaic]\label{lem:V11-implies-M3}
If $\Q$ satisfies (V11) and every mosaic $M\in\Mos_\Q$ is
$L_\Q$-compatible in the sense of Definition~\ref{def:LM-Lq}, then
$M$ satisfies the round-4 mosaic axiom (M3).
\end{lemma}

\begin{proof}
For $p,q,r\in P$ at positions $\pi_p,\pi_q,\pi_r\in\Pos_\Q$, the
mosaic labels are $L_M(p,q)=L_\Q(\pi_p,\pi_q)$ (and similarly for
the other two pairs).  By (V11), $L_\Q(\pi_p,\pi_r)\in\comp(L_\Q(\pi_p,\pi_q),
L_\Q(\pi_q,\pi_r))$, so $L_M(p,r)\in\comp(L_M(p,q),L_M(q,r))$,
which is (M3).
\end{proof}

\subsection{Restated support-closure}

The round-4 support-closure clauses (SC1)--(SC4) are restated with
explicit reference to $L_\Q$ and (V11) replacing the implicit
composition-on-$\T_\Q$.

\begin{definition}[Support-closed mosaic family, non-circular]\label{def:support-closed-v3}
A finite family $\Mos$ of mosaics is support-closed relative to
$(\T_\Q,L_\Q)$ if it satisfies (SC1)--(SC3) of \cite{GPT55Repaired}
unchanged, together with:
\begin{enumerate}[label=(SC4$'$)]
\item For every abstract position symbol $\pi\in\Pos_\Q$ on the
generated context of some source in $\Ctx_\Q$, $\Mos$ contains a
mosaic $M_\pi$ that includes $\pi$ as one of its positions, with
$\tau_{M_\pi}(\pi)=\lambda(\pi)$ (the type declared by the certificate)
and the stratum of $\pi$ in $M_\pi$ matching its declared stratum.
\end{enumerate}
\end{definition}

The change is that (SC4$'$) no longer requires per-triple coverage
of $\Trip_\Q$.  Instead, the composition table on $\T_\Q$ is
enforced \emph{globally} by (V11), and (SC4$'$) requires only that
each position has \emph{some} mosaic realising its type and
stratum.  Triple consistency is then a consequence of (V11) and
(SC1)--(SC3), without circular appeal to per-triple mosaic
coverage.

\subsection{Restated patchwork theorem}

\begin{theorem}[Support-closed mosaic patchwork, non-circular]\label{thm:patchwork-v3}
Let $\Q$ be a valid certificate satisfying (V1)--(V13), (M6$'$) on
every mosaic in $\Mos_\Q$, (V11) on the abstract pair-label function
$L_\Q$, and (SC1)--(SC3), (SC4$'$).  Then every finite prefix
$F\subseteq\U(\Q)$ satisfies the abstract RCC5 frame axioms
(F1)--(F3) on all triples of occurrences.
\end{theorem}

\begin{proof}
We argue independently of the mosaic family on the global
relation:

\emph{Step 1.}  Define the unfolding label $\lab(u,v)\in\Base$ as
$L_\Q(\pi_u,\pi_v)$, where $\pi_u,\pi_v$ are the abstract position
symbols of the unfolding occurrences $u,v$ in $\U(\Q)$ (each
occurrence carries a unique abstract position symbol by the
unfolding rule of \S4 of \cite{GPT55Repaired}, possibly via an
equality-port chain).

\emph{Step 2.}  (F1) and (F2) follow from (L1) and (L2) of
Definition~\ref{def:Lq}.  For (F3):  given any triple $(u,v,w)$ in
$F$, with abstract position symbols
$(\pi_u,\pi_v,\pi_w)$, apply (V11) directly:
$\lab(u,w)=L_\Q(\pi_u,\pi_w)\in\comp(L_\Q(\pi_u,\pi_v),L_\Q(\pi_v,\pi_w))
=\comp(\lab(u,v),\lab(v,w))$.  No mosaic is consulted; the
composition check is a finite global validity check on $L_\Q$.

\emph{Step 3.}  The mosaic family $\Mos_\Q$ is used only to discharge
existential support obligations (M4)--(M5) and the
$\EQ$-congruence/universal-propagation conditions, which the
round-4 lemmas establish on the mosaic side (Lemmas~5.7, 5.8, 5.9 of
\cite{GPT55Repaired}).  These do not rely on (SC4) per-triple
coverage; they are local to a single mosaic.  Mosaic
type-consistency at each abstract position is supplied by (SC4$'$).
\end{proof}

\begin{remark}
The repair shifts the burden of composition consistency from
mosaic-mediated (SC4) coverage to a global validity clause on the
finite $L_\Q$.  The proof of (F3) no longer requires the existence
of a mosaic per triple, and no statement of the patchwork theorem
assumes its own conclusion.
\end{remark}

%% ==================================================================
\section{Item 5: strengthened overlap amalgamation}
\label{sec:item5}
%% ==================================================================

\subsection{The weakness diagnosed by round 5}

Round-4's overlap-amalgamation lemma (Lemma~5.6 of
\cite{GPT55Repaired}) takes hypothesis~(O3):
\begin{quote}
every triple $(u,v,w)$ with $u\in P_1\setminus P_0$,
$w\in P_2\setminus P_0$, $v\in P_0$ admits some
$\ell\in\Base$ with $\ell\in\comp(L_1(u,v),L_2(v,w))$.
\end{quote}
The construction then takes $L_{12}(u,w)\in\bigcap_{v\in P_0}
\comp(L_1(u,v),L_2(v,w))$.  Round 5~\cite{GPT55Round5Review} notes
that per-triple existence of some $\ell$ does not imply that the
intersection is non-empty, nor that one consistent assignment
exists across all cross-pairs simultaneously.

\subsection{Strengthened hypothesis (O3$'$) and validity clause (V12)}

\begin{definition}[Strong overlap compatibility]\label{def:strong-overlap}
Two mosaics $M_1,M_2\in\Mos_\Q$ are \emph{strongly
overlap-compatible} over $P_0=P_1\cap P_2$ if there exists a
declared cross-label function
\[
  L_{12}^{\rm cross}:(P_1\setminus P_0)\times(P_2\setminus P_0)
                     \;\longrightarrow\;\Base
\]
satisfying:
\begin{enumerate}[label=(O3$'$.\alph*)]
\item \textbf{Triangle compatibility through every overlap node.}
For every $u\in P_1\setminus P_0$, $w\in P_2\setminus P_0$, and
every $v\in P_0$,
\[
  L_{12}^{\rm cross}(u,w)\;\in\;\comp\bigl(L_1(u,v),L_2(v,w)\bigr).
\]
\item \textbf{Composition consistency on mixed triples.}  For every
$u,u'\in P_1\setminus P_0$, $w\in P_2\setminus P_0$,
\[
  L_{12}^{\rm cross}(u',w)\;\in\;
    \comp\bigl(L_1(u',u),L_{12}^{\rm cross}(u,w)\bigr),
\]
and symmetrically for $u\in P_1\setminus P_0$,
$w,w'\in P_2\setminus P_0$.
\item \textbf{Inverse symmetry.}  $L_{12}^{\rm cross}(w,u)=
\Inv\bigl(L_{12}^{\rm cross}(u,w)\bigr)$.
\end{enumerate}
\end{definition}

\paragraph{(V12) Declared overlap labels.}
For every pair of mosaics $M_1,M_2\in\Mos_\Q$ that share a non-empty
boundary $P_0=P_1\cap P_2$ and whose amalgamation is required by
(SC3), the certificate declares a cross-label function
$L_{12}^{\rm cross}$ satisfying (O3$'$.a)--(O3$'$.c) of
Definition~\ref{def:strong-overlap}.  In particular, the declared
$L_{12}^{\rm cross}$ extends consistently to a global mosaic
$M_{12}=(P_1\cup P_2,\tau_{12},L_{12})$ with $L_{12}\subseteq L_\Q$
(consistent with Definition~\ref{def:LM-Lq}).

The clause (V12) is finite-checkable: the certificate enumerates
the pairs of mosaics requiring amalgamation (finite, since
$\Mos_\Q$ is finite), and the cross-label function for each pair is
a finite map.  The composition-table checks of (O3$'$.a)--(O3$'$.c)
are finite quantifiers over $|P_1\setminus P_0|\cdot|P_0|\cdot
|P_2\setminus P_0|$ triples, each a constant-time lookup.

\subsection{Restated overlap amalgamation lemma}

\begin{lemma}[Overlap amalgamation, strong]\label{lem:overlap-v3}
Let $M_1,M_2\in\Mos_\Q$ be strongly overlap-compatible
(Definition~\ref{def:strong-overlap}), and let $L_{12}^{\rm cross}$
be the declared cross-label function from (V12).  Then the
amalgamated mosaic
$M_{12}=(P_1\cup P_2,\tau_{12},L_{12},\mathcal V_{12})$ with
\[
  L_{12}(p,q):=\begin{cases}
    L_i(p,q) & \text{if }p,q\in P_i,\\
    L_{12}^{\rm cross}(p,q) & \text{if }p\in P_1\setminus P_0,\;q\in P_2\setminus P_0,\\
    \Inv L_{12}^{\rm cross}(q,p) & \text{if }p\in P_2\setminus P_0,\;q\in P_1\setminus P_0,
  \end{cases}
\]
satisfies (M1)--(M5) and (M6$'$) of v2~\cite{V2Delta} on all
positions of $P_1\cup P_2$.
\end{lemma}

\begin{proof}
(M1) and (M2) follow from (O3$'$.c) and the unchanged structure on
$P_1, P_2$.  (M3) on a triple $(u,v,w)$ in $P_{12}$:\@
\begin{itemize}
\item If $u,v,w$ all lie in $P_i$ for some $i$, (M3) on $M_i$
suffices.
\item If two of $u,v,w$ lie in $P_i$ and the third in
$P_j\setminus P_0$ ($i\neq j$), apply (O3$'$.b) to the cross pair
through the third overlap vertex (or one of the within-$P_i$
vertices), which is a triangle of the form covered by (O3$'$.a)
or (O3$'$.b).
\item If $u\in P_1\setminus P_0$, $v\in P_0$, $w\in P_2\setminus
P_0$, (O3$'$.a) directly gives
$L_{12}^{\rm cross}(u,w)\in\comp(L_1(u,v),L_2(v,w))$.
\item All other cases reduce to (O3$'$.b) by symmetry.
\end{itemize}
(M4) and (M5) are inherited from $M_1, M_2$ as in round-4
Lemmas~5.7 and 5.8.  (M6$'$) is preserved: (M6$'$.b) requires
co-vertical $\PP/\PPI$ labels;\@ cross-pair labels chosen by
$L_{12}^{\rm cross}$ are in $\{\DR,\PO\}$ on cross-stratum pairs
because (O3$'$) does not produce $\PP/\PPI$ across distinct strata.
(M6$'$.c)--(M6$'$.d) are similarly preserved.
\end{proof}

\begin{remark}
Round-4's existential ``some $\ell\in\Base$'' in (O3) is replaced
by an explicit declared function $L_{12}^{\rm cross}$.  The
ambiguity ``does the intersection over $v\in P_0$ have a common
witness?'' is moved to the certificate-writing time:\@ the writer
must produce a witness function and (V12) checks it.
\end{remark}

%% ==================================================================
\section{Item 6: joint equality validity clause}
\label{sec:item6}
%% ==================================================================

\subsection{The diagnosis}

Round-4 \cite{GPT55Repaired} proves the standalone typed-$\EQ$
quotient theorem~7.3 under six assumptions (A1)--(A6).  The
delicate one is (A4):
\begin{quote}
For all $u\eqc v$ and every $\exists R.D\in\lambda(u)=\lambda(v)$
with $R\in\{\PP,\PPI\}$, there exists $w\in\Omega$ with
$\{u,v\}\Reach_R w$ (equality-aware reachability) and
$D\in\lambda(w)$.  Equivalently, a single witness discharges the
obligation simultaneously at every $\eqc$-mate of the source.
\end{quote}
Round-4 derives (A4) from (V4) + the definition of $\Reach$.
Round-5 \cite{GPT55Round5Review} suggests promoting (A4) to an
explicit finite validity clause.

\subsection{The repair: clause (V13)}

\paragraph{(V13) Joint equality discharge.}
For every mosaic $M\in\Mos_\Q$, every stratum $V\in\Stratum_M$,
every explicit equality class $E\subseteq P_V$ (entirely within $V$
by (M6$'$.a)), and every $\exists R.D\in\Cl(C_0)$ such that
$\exists R.D\in\tau(p)$ for all $p\in E$, the certificate declares a
joint witness $w_E^{R,D}\in P_V$ with $L_M(p,w_E^{R,D})=R$ and
$D\in\tau(w_E^{R,D})$ for some (equivalently, every) $p\in E$.

In addition, for cross-class joint witnesses, when an equality
class $E\subseteq P_V$ has $\exists R.D$ in the common type and
the existing port-indexed mate cluster $\Mate_\Q(C_V,s_V,p)$ for
$p\in E$ contains a port $(s',p')$ targeting a witness in
another generated context, the certificate declares that
$\Mate_\Q(C_V,s_V,p)$ is identical for every $p\in E$ (the cluster
is class-wide, not per-port-of-the-class).

\subsection{(V13) implies (A4)}

\begin{lemma}[(V13) discharges (A4)]\label{lem:V13-implies-A4}
Suppose $\Q$ satisfies (V1)--(V13).  Then the lift of $\eqc^Q$ to
$\U(\Q)$ satisfies assumption~(A4) of the round-4 typed-$\EQ$
quotient theorem.
\end{lemma}

\begin{proof}
Take $u,v\in\Omega(\U(\Q))$ with $u\eqc v$ and
$\exists R.D\in\lambda(u)=\lambda(v)$ with $R\in\{\PP,\PPI\}$.  By
the unfolding rule, $u$ and $v$ realise positions $p_u,p_v$ of some
mosaic $M\in\Mos_\Q$ with $p_u,p_v$ in the same explicit equality
class $E$ of $M$.  By (M6$'$.a), $E\subseteq P_V$ for some
stratum $V$.  Since $u\eqc v$ at the certificate level, both
$p_u$ and $p_v$ belong to the same class $E$.

By (V13), $\Q$ declares a joint witness $w_E^{R,D}\in P_V$ with
$L_M(p,w_E^{R,D})=R$ and $D\in\tau(w_E^{R,D})$ for every $p\in E$.
Take $w$ to be the occurrence in $\U(\Q)$ realising
$w_E^{R,D}$.  Then $L_M(p_u,w_E^{R,D})=R$ and
$L_M(p_v,w_E^{R,D})=R$ both hold, so $u\Reach_R w$ and
$v\Reach_R w$ via the same witness occurrence.  This is exactly
(A4).
\end{proof}

\begin{remark}
The strengthening eliminates the round-4 inference chain through
the equality-aware reachability $\Reach$.  Where round 4 used the
\emph{existence} of a witness for each individual $\eqc$-mate, (V13)
declares the witness once for the entire class.  This is the
formal counterpart of the round-5 review's recommended repair:
``promote (A4) to an explicit finite validity clause''.
\end{remark}

%% ==================================================================
\section{Item 7: half-open cycle convention}
\label{sec:item7}
%% ==================================================================

\subsection{The off-by-one risk in the round-4 convention}

Round-4 \cite{GPT55Repaired} states the request-fulfilled
repetition lemma (Lemma~6.12) with the pumped path
\[
  u_0,u_1,\ldots,u_{i-1},(u_i,\ldots,u_{j-1})^\omega.
\]
The cycle word is $u_i,u_{i+1},\ldots,u_{j-1}$; position $u_j$ is
not part of the cycle word, but acts as the ``return-to-$u_i$'' on
the next lap.  Round-5 \cite{GPT55Round5Review} notes that
round-4 alternates between a closed segment $u_i,\ldots,u_j$ and a
lasso that repeats $u_i,\ldots,u_{j-1}$, which is the standard
off-by-one source of confusion.  The lemma is correct;\@ the
convention should be made explicit.

\subsection{The repair: half-open $[i,j)$ convention with explicit endpoint mapping}

\begin{definition}[Half-open cycle convention]\label{def:half-open}
The cycle word of a request-fulfilled repetition lemma is the
\emph{half-open segment} $u_i,u_{i+1},\ldots,u_{j-1}$ of length
$j-i$.  The pumped path is the concatenation of $u_0,\ldots,u_{i-1}$
with infinitely many fresh copies of the cycle word:
\[
  u_0,u_1,\ldots,u_{i-1},
  \underbrace{u_i^{(1)},\ldots,u_{j-1}^{(1)}}_{\text{lap }1},
  \underbrace{u_i^{(2)},\ldots,u_{j-1}^{(2)}}_{\text{lap }2},
  \ldots
\]
For every $\ell\geq 1$, the occurrence $u_j^{(\ell)}$ \emph{does
not appear} in the pumped path; the next position after
$u_{j-1}^{(\ell)}$ is $u_i^{(\ell+1)}$.  At the profile level,
$u_j^{(\ell)}$ is identified with $u_i^{(\ell+1)}$ because
$\Pi(u_j)=\Pi(u_i)$ by the pigeonhole choice of $i<j$.
\end{definition}

\begin{lemma}[Request-fulfilled repetition lemma, half-open]\label{lem:request-cycle-v3}
Let $u_0,u_1,u_2,\ldots$ be an infinite generated vertical path in
a split presentation, and suppose every transitive existential
request on the path is eventually fulfilled.  Then there are
indices $i<j$ with $\Pi(u_i)=\Pi(u_j)$ and the half-open segment
$[u_i,u_j)$ is \emph{request-fulfilled}: for every transitive
existential request $\exists R.D$ pending at some position $u_k$
with $i\leq k<j$, there is a strict later position $u_{k'}$ within
$[u_i,u_j)$ (i.e., $k<k'<j$) or in the segment endpoint $u_j$
(handled below) with $D\in\lambda(u_{k'})$.

The pumped path
$u_0,\ldots,u_{i-1},u_i^{(1)},\ldots,u_{j-1}^{(1)},
u_i^{(2)},\ldots,u_{j-1}^{(2)},\ldots$ built from fresh
occurrences of the half-open segment satisfies the strong
fulfillment property:
\begin{quote}
every $\PP/\PPI$ existential request generated at any occurrence
of the pumped path is fulfilled at a strict later occurrence on
the same path, within bounded lookahead $\leq j-i$.
\end{quote}
\end{lemma}

\begin{proof}
The pigeonhole and checkpoint arguments are unchanged from
round-4 (Lemma~6.10 of \cite{GPT55Repaired}).  We make the endpoint
case explicit.

\emph{Endpoint case.}  Let $\exists R.D$ be a request pending at
$u_k$ for some $i\leq k<j$.  By round-4's analysis~(F1) and
(F2a)/(F2b):
\begin{enumerate}[label=(\alph*)]
\item if $\exists R.D$ is fulfilled at some $u_{k'}$ with
$k<k'<j$, the pumped path contains $u_{k'}^{(\ell)}$ strictly
later than $u_k^{(\ell)}$ in the same lap.  Lookahead:
$k'-k<j-i$.

\item if $\exists R.D$ is fulfilled at $u_j$ in the original path,
the pumped path's next-lap position $u_i^{(\ell+1)}$ has
$\Pi(u_i^{(\ell+1)})=\Pi(u_j)$ at the profile level, but at the
\emph{semantic} level $u_i^{(\ell+1)}$ is a fresh occurrence with
$\lambda(u_i^{(\ell+1)})=\lambda(u_i)$, not $\lambda(u_j)$.  In
general $\lambda(u_i)$ and $\lambda(u_j)$ are equal as types
(since profiles include closure types in round-4 Definition~6.6 of
\cite{GPT55Repaired}), so $D\in\lambda(u_i^{(\ell+1)})$ iff
$D\in\lambda(u_j)$.  Therefore the pumped path's $u_i^{(\ell+1)}$
acts as the strict later fulfilling position for $\exists R.D$
generated at $u_k^{(\ell)}$.  Lookahead: $(j-k)+0=j-k\leq j-i$.

\item if $\exists R.D$ persists into $\Pend(u_j)=\Pend(u_i)$,
case~(F2b) of round-4 applies: there is $m$ with $i\leq m<j$ and
$D\in\lambda(u_m)$;\@ the pumped path's $u_m^{(\ell+1)}$ in the
next lap satisfies $D\in\lambda(u_m^{(\ell+1)})$, strictly later
than $u_k^{(\ell)}$.  Lookahead: $(j-k)+(m-i)+1\leq 2(j-i)$.
\end{enumerate}
The half-open convention makes case~(b) explicit: $u_j$ never
appears in the pumped path, but its \emph{profile} is realised by
the next lap's $u_i^{(\ell+1)}$, and (since profiles include
closure types) the witness label transfers to the next lap.
\end{proof}

\begin{remark}
The change is purely notational: round-4's lemma is correct, but
the closed-vs-half-open ambiguity is removed by always using
$[i,j)$ and stating the endpoint behaviour explicitly.  The
boundary $u_j$ is never an occurrence of the pumped path; the
position right after the last cycle-word occurrence of a lap is
the first cycle-word occurrence of the next lap.
\end{remark}

%% ==================================================================
\section{Item 8: global finite regularization}
\label{sec:item8}
%% ==================================================================

\subsection{The diagnosis}

Round-4 \cite{GPT55Repaired} proves the request-fulfilled
repetition lemma per-ray (Lemma~6.12 and Lemma~6.13 of
\cite{GPT55Repaired}): each infinite vertical ray of a split
presentation independently lassos to a finite cycle.  Round-5
\cite{GPT55Round5Review} requests a global lemma showing that a
single finite certificate covers \emph{all} infinite vertical rays
of the branching split forest uniformly, not infinitely many
independent lasso choices.

\subsection{The repair: a global representative-selection lemma}

The key observation is that the finite extended-profile alphabet
$\extSigma$ is finite of size at most $|\Sigma|\cdot 2^n$, and a
lasso choice depends only on the extended profile $\extProf$ at the
two pigeonhole positions, not on the specific ray.

\begin{lemma}[Uniform representative selection]\label{lem:global-rep}
Let $\extSigma$ be the finite extended-profile alphabet (round-4
Definition~6.7).  For every $\sigma\in\extSigma$, choose:
\begin{enumerate}[label=(\arabic*)]
\item a representative finite generated context
$C_\sigma\in\Ctx_\Q$ whose source has extended profile $\sigma$;
\item a representative cycle word
$\mathsf{cyc}_\sigma=(u_0^\sigma,u_1^\sigma,\ldots,u_{L_\sigma-1}^\sigma)$
of length $L_\sigma$, built from $C_\sigma$ using the construction
of Lemma~\ref{lem:request-cycle-v3}, such that the half-open
segment $\mathsf{cyc}_\sigma$ is request-fulfilled.
\end{enumerate}
If, for every $\sigma\in\extSigma$, such a representative
$(C_\sigma,\mathsf{cyc}_\sigma)$ exists, then a single finite
certificate $\Q$ covers every infinite vertical ray of the split
forest, by uniformly substituting at every ray the canonical
$\mathsf{cyc}_{\sigma(r)}$ for the extended profile $\sigma(r)$
that the ray hits twice.
\end{lemma}

\begin{proof}
We argue in two parts.

\emph{Part 1 (Existence of representatives).}  By the round-4
hypothesis that every transitive existential request is eventually
fulfilled in the witness-generated submodel, every infinite
vertical ray $r$ satisfies the hypothesis of Lemma~6.10 (checkpoint
structure) of \cite{GPT55Repaired}.  Hence for every $r$ the
pigeonhole gives indices $i_r<j_r$ with
$\extProf(u_{i_r}^r)=\extProf(u_{j_r}^r)=\sigma(r)$.

For each $\sigma\in\extSigma$ realised on some ray, fix one such
ray $r_\sigma$ and set $C_\sigma$ to be the generated context of
$u_{i_{r_\sigma}}^{r_\sigma}$ and
$\mathsf{cyc}_\sigma=[u_{i_{r_\sigma}}^{r_\sigma},\ldots,
u_{j_{r_\sigma}-1}^{r_\sigma})$.  This is the
\emph{Axiom-of-Choice-free} representative selection:\@ for each of
the finitely many extended profiles in $\extSigma$, choose any
one ray realising it.  Since $\extSigma$ is finite, the selection
is finite and effective.

\emph{Part 2 (Uniform substitution).}  Consider any infinite
vertical ray $r$ of the split forest.  By Part~1, $r$ hits some
extended profile $\sigma(r)\in\extSigma$ at two positions $i_r<j_r$
with $\extProf(u_{i_r}^r)=\extProf(u_{j_r}^r)$.  By the
substitution lemma of \cite{GPT55Repaired} (Lemma~6.9),
substituting the representative cycle $\mathsf{cyc}_{\sigma(r)}$
for the tail $[u_{i_r}^r,\ldots]$ of $r$ preserves all closure
formulas at every position of $r$, because both
$(u_{i_r}^r,\ldots)$ and $\mathsf{cyc}_{\sigma(r)}$ realise the
same extended profile at their starting positions and are
request-fulfilled.

The substitution is the \emph{same} cycle word for every ray that
hits $\sigma(r)$;\@ rays that hit different extended profiles
substitute different (but pre-declared) cycle words.  The total
substitution data is a finite map
$\extSigma\to\{(C_\sigma,\mathsf{cyc}_\sigma):\sigma\in\extSigma\}$,
finite of size $\leq |\extSigma|$.

After substitution, every infinite vertical ray of the (modified)
split forest is the original prefix
$u_0^r,\ldots,u_{i_r-1}^r$ concatenated with
$\mathsf{cyc}_{\sigma(r)}^\omega$.  The finitely many cycle words
$\{\mathsf{cyc}_\sigma\}_{\sigma\in\extSigma}$ together with the
finite prefix data of the unmodified portion form a finite
certificate.
\end{proof}

\paragraph{(V14) Uniform cycle representatives.}
For each $\sigma\in\extSigma$ that is realised on some infinite
vertical ray of the witness-generated submodel, the certificate
$\Q$ declares one representative finite generated context $C_\sigma$
and one representative half-open cycle word $\mathsf{cyc}_\sigma$
satisfying the request-fulfilled property of
Lemma~\ref{lem:request-cycle-v3}.  Uniformity:\@ no two distinct
extended profiles use the same cycle word, and the substitution
function is the identity on the finite prefix
$u_0,\ldots,u_{i-1}$ before the first pigeonhole hit.

(V14) is a finite quantifier over $|\extSigma|\leq|\Sigma|\cdot
2^n$ extended profiles;\@ for each, the certificate declares a
context (one of $|\Ctx_\Q|$ many) and a cycle word (a path in
$\U(\Q)$ of length at most $|\Ctx_\Q|\cdot|S|$).

\begin{corollary}[Global finite regularization]\label{cor:global-reg}
A valid certificate $\Q$ satisfying (V1)--(V14) yields, by the
uniform substitution of Lemma~\ref{lem:global-rep}, a single finite
representation of every infinite vertical ray of the original
witness-generated submodel.  No automaton, parity acceptance, or
$\omega$-language theoretic machinery is required;\@ the
pigeonhole is over the finite set $\extSigma$.
\end{corollary}

%% ==================================================================
\section{Summary and combined statement of the proof}
\label{sec:summary}
%% ==================================================================

The round-5 review~\cite{GPT55Round5Review} lists eight
recommended repairs.  Combining this paper with the v2 (M6$'$)
delta~\cite{V2Delta} closes all eight:

\begin{itemize}
\item \textbf{Item~1 (TeX/PDF mismatch).}  Self-consistent in the
repository (the source compiles to the 31-page PDF; the reviewer's
bundle had stale source).  Editorial; no mathematical change.
\item \textbf{Items~2, 3, 9.}  Closed by~\cite{V2Delta}: recursive
verticalization (M6$'$), the stress concept
$\exists\DR.(A\sqcap\exists\PP.B)\sqcap\exists\DR.B$ accepted by
the worked example, polynomial side-width $m(n)=O(n^3)$ replaced
by an effectively computable recurrence over modal depth.
\item \textbf{Item~4.}  Closed by \S\ref{sec:item4}: explicit
abstract pair-label function $L_\Q$ and validity clause (V11) for
the composition table on $\T_\Q$.  The patchwork theorem is
non-circular.
\item \textbf{Item~5.}  Closed by \S\ref{sec:item5}: strong overlap
compatibility (O3$'$) and validity clause (V12) requiring a
declared global cross-label assignment.
\item \textbf{Item~6.}  Closed by \S\ref{sec:item6}: (V13) promotes
the joint equality discharge from assumption (A4) to an explicit
finite validity clause.
\item \textbf{Item~7.}  Closed by \S\ref{sec:item7}: half-open
$[i,j)$ convention with explicit endpoint mapping
$u_j\mapsto u_i^{(\ell+1)}$ at the profile level.
\item \textbf{Item~8.}  Closed by \S\ref{sec:item8}: uniform
representative selection (V14) gives a single finite certificate
covering all infinite vertical rays.
\end{itemize}

\subsection{Combined validity-clause numbering}

The combined certificate validity check has fourteen clauses:
(V1)--(V10) of round-4~\cite{GPT55Repaired}, modified in v2 to
(V3$'$), (V4$'$), (V5$'$), (V9.Eb$'$)~\cite{V2Delta}, together with
(V11)--(V14) of this paper.  Each clause is a finite quantifier
over the certificate's finite alphabets and contexts;\@ no clause
requires a quantifier over models or unfoldings.

\subsection{Combined decision procedure}

The decision procedure for $\ALCIRCC{5}$ concept satisfiability is:
\begin{enumerate}[label=(\arabic*)]
\item enumerate finite generated split-forest certificates $\Q$ of
size $\leq B(C_0)\leq 2^{2^{p(n)}}$ for an explicit computable
polynomial $p$, where the components of $\Q$ are:
$S$ (profiles), $E_\Q$ (equality ports), $\Ctx_\Q$ (context
schemes), $\Mos_\Q$ (mosaics with stratum forests $\mathcal V$),
$L_\Q$ (abstract pair-label function),
$\{L_{12}^{\rm cross}\}$ (declared overlap labels),
$\{w_E^{R,D}\}$ (joint equality witnesses),
$\{(C_\sigma,\mathsf{cyc}_\sigma)\}_{\sigma\in\extSigma}$ (uniform
cycle representatives);
\item for each $\Q$, check (V1)--(V14) and (M6$'$.a)--(M6$'$.d), all
finite-graph computations on $\Q$'s alphabets;\@
\item if any $\Q$ passes, output ``$C_0$ is satisfiable'';
otherwise output ``$C_0$ is unsatisfiable''.
\end{enumerate}
Soundness is the v2 vertical-realisation
lemma~\cite{V2Delta} combined with Theorem~\ref{thm:patchwork-v3}
and the typed-$\EQ$ quotient theorem of round 4 with (A4) replaced
by (V13).  Completeness is the uniform regularization
Corollary~\ref{cor:global-reg} combined with the existence of
representatives from witness-generated submodels of any
satisfying abstract RCC5 frame, as in round 4.

%% ==================================================================
\begin{thebibliography}{99}

\bibitem{GPT55Repaired}
GPT-5.5 (OpenAI), prompted by Michael Wessel.
\emph{A Generated Split-Forest Decidability Proof for $\ALCIRCC{5}$
Without Automata: A Repaired Containment-Oriented Presentation}.
Unpublished manuscript, May 2026.  Repository:
\url{https://github.com/lambdamikel/alcircc5/blob/master/papers/gpt5.5_round2/repaired_split_forest_no_automata_proof.pdf}.

\bibitem{GPT55Round5Review}
GPT-5.5 (OpenAI), prompted by Michael Wessel.
\emph{Formal Review of the Latest Repaired Split-Forest Proof
Manuscript}.  Unpublished manuscript, May~22, 2026.  Repository:
\url{https://github.com/lambdamikel/alcircc5/blob/master/papers/gpt5.5_round2/formal_review_progress_paper.pdf}.

\bibitem{V2Delta}
Claude (Anthropic, Opus~4.7), prompted by Michael Wessel.
\emph{Repaired Split-Forest Decidability Proof for $\ALCIRCC{5}$ (v2):
Recursive Verticalization (M6$'$) to Close the Round-5
Side-Witness Gap}.  Unpublished manuscript, May~2026.  Repository:
\url{https://github.com/lambdamikel/alcircc5/blob/master/papers/gpt5.5_round2/repaired_split_forest_no_automata_proof_v2.pdf}.

\bibitem{RenzNebel1999}
Jochen Renz and Bernhard Nebel.
\emph{On the Complexity of Qualitative Spatial Reasoning: A
Maximal Tractable Fragment of the Region Connection Calculus}.
Artificial Intelligence, 108(1--2):69--123, 1999.

\bibitem{ALCIRCC5Repo}
Michael Wessel et al.
\emph{ALCI\_RCC5 Project Repository}.
\url{https://github.com/lambdamikel/alcircc5}.

\end{thebibliography}

\end{document}
